% Part: first-order-logic
% Chapter: tableaux
% Section: provability-consistency

\documentclass[../../../include/open-logic-section]{subfiles}

\begin{document}

\iftag{FOL}
      {\olfileid[fr]{fol}{tab}{prv}}
      {\olfileid[fr]{pl}{tab}{prv}}

\olsection{\usetoken{S}{derivability} et cohérence}

Établissons plusieurs propriétés de la relation de !!{derivability}.
Outre leur intérêt propre, elles interviendront toutes dans la
démonstration du théorème de complétude.

\begin{prop}\ollabel{prop:provability-contr}
  Si $\Gamma \Proves !A$ et si $\Gamma \cup \{!A\}$ est
  incohérent, alors $\Gamma$ est incohérent.
\end{prop}

\begin{proof}
Il existe deux parties finies de~$\Gamma$,
$\Gamma_0 = \{!B_1, \dots, !B_n\}$ et
$\Gamma_1 = \{!C_1, \dots, !C_m\}$, telles que les deux ensembles
  \begin{align*}
    \{\sFmla{\False}{!A}, &
    \sFmla{\True}{!B_1}, \dots, \sFmla{\True}{!B_n}\} \\
    \{\sFmla{\True}{!A}, &
    \sFmla{\True}{!C_1}, \dots, \sFmla{\True}{!C_m}\}
  \end{align*}
possèdent des !!{tableau}s fermés. En appliquant la règle \Cut{} à
$!A$, nous pouvons les combiner en un seul !!{tableau} fermé, qui
montre que $\Gamma_0 \cup \Gamma_1$ est incohérent. Comme
$\Gamma_0 \subseteq \Gamma$ et $\Gamma_1 \subseteq \Gamma$, on a
$\Gamma_0 \cup \Gamma_1 \subseteq \Gamma$; par conséquent,
$\Gamma$ est incohérent.
\end{proof}

\begin{editorial}
Dans cette preuve, la source définit $\Gamma_1$ au moyen de
$!C_1,\dots,!C_n$, puis emploie $!C_m$. L'indice~$m$ est rétabli dans
la définition de $\Gamma_1$.
\end{editorial}

\begin{prop}
\ollabel{prop:prov-incons}
$\Gamma \Proves !A$ si et seulement si
$\Gamma \cup \{\lnot !A\}$ est incohérent.
\end{prop}

\begin{proof}
Supposons d'abord que $\Gamma \Proves !A$. Il existe donc des
$!B_1,\dots,!B_n \in \Gamma$ et un !!{tableau} fermé pour
\[
\{\sFmla{\False}{!A},
\sFmla{\True}{!B_1}, \dots, \sFmla{\True}{!B_n}\}
\]
Remplaçons l'hypothèse $\sFmla{\False}{!A}$ par
$\sFmla{\True}{\lnot !A}$ et appliquons la règle
$\TRule{\True}{\lnot}$, dont la conclusion est précisément
$\sFmla{\False}{!A}$. On obtient ainsi un !!{tableau} fermé pour
\[
\{\sFmla{\True}{\lnot !A},
\sFmla{\True}{!B_1}, \dots, \sFmla{\True}{!B_n}\}.
\]

Réciproquement, supposons qu'il existe un !!{tableau} fermé pour ce
dernier ensemble. Supprimons l'hypothèse
$\sFmla{\True}{\lnot !A}$ ainsi que toute formule obtenue en lui
appliquant $\TRule{\True}{\lnot}$, puis ajoutons l'hypothèse
$\sFmla{\False}{!A}$. Si une branche se fermait parce qu'elle contenait
la conclusion de $\TRule{\True}{\lnot}$ appliquée à
$\sFmla{\True}{\lnot !A}$, c'est-à-dire
$\sFmla{\False}{!A}$, la branche correspondante du nouveau
!!{tableau} est encore fermée. Si une branche de l'ancien tableau se
fermait parce qu'elle contenait à la fois l'hypothèse
$\sFmla{\True}{\lnot !A}$ et $\sFmla{\False}{\lnot !A}$, appliquons
$\TRule{\False}{\lnot}$ à cette dernière formule. Nous obtenons
$\sFmla{\True}{!A}$, qui ferme la branche avec la nouvelle hypothèse
$\sFmla{\False}{!A}$.
\end{proof}

\begin{prob}
Démontrer que $\Gamma \Proves \lnot !A$ si et seulement si
$\Gamma \cup \{!A\}$ est incohérent.
\end{prob}

\begin{prop}\ollabel{prop:explicit-inc}
  Si $\Gamma \Proves !A$ et $\lnot !A \in \Gamma$, alors
  $\Gamma$ est incohérent.
\end{prop}

\begin{proof}
  Supposons que $\Gamma \Proves !A$ et $\lnot !A \in \Gamma$.
  Il existe alors $!B_1, \dots, !B_n \in \Gamma$ tels que
  \[
  \{\sFmla{\False}{!A}, \sFmla{\True}{!B_1}, \dots, \sFmla{\True}{!B_n}\}
  \]
  possède un !!{tableau} fermé. Remplaçons l'hypothèse
  \sFmla{\False}{!A} par \sFmla{\True}{\lnot !A}, puis insérons,
  après les $n+1$ hypothèses, la conclusion de
  \TRule{\True}{\lnot} appliquée à \sFmla{\True}{\lnot !A}, à savoir
  \sFmla{\False}{!A}. Tout !!{sentence} du !!{tableau} dont la
  justification renvoyait à la ligne~$1$ de l'ancien !!{tableau}
  renvoie maintenant à la ligne~$n+2$. Si l'ancien !!{tableau} était
  fermé, le nouveau l'est donc aussi. Comme toutes ses hypothèses
  appartiennent à~$\Gamma$, il montre que $\Gamma$ est incohérent.
\end{proof}

\begin{editorial}
La source applique ici par erreur $\TRule{\True}{\lnot}$ à
$\sFmla{\False}{!A}$ et renvoie ensuite à la ligne~$n+1$. La règle
s'applique à $\sFmla{\True}{\lnot !A}$; sa conclusion, placée après
les $n+1$ hypothèses, occupe la ligne~$n+2$.
\end{editorial}

\begin{prop}\ollabel{prop:provability-exhaustive}
  Si $\Gamma \cup \{!A\}$ et $\Gamma \cup \{\lnot !A\}$ sont tous
  deux incohérents, alors $\Gamma$ est incohérent.
\end{prop}

\begin{proof}
  Supposons qu'il existe $!B_1, \dots, !B_n \in \Gamma$ et
  $!C_1, \dots, !C_m \in \Gamma$ tels que les deux ensembles
  \begin{align*}
    \{\sFmla{\True}{!A}, &
    \sFmla{\True}{!B_1}, \dots, \sFmla{\True}{!B_n}\} \text{ et}\\
    \{\sFmla{\True}{\lnot !A}, &
    \sFmla{\True}{!C_1}, \dots, \sFmla{\True}{!C_m}\}
  \end{align*}
  possèdent chacun un !!{tableau} fermé. Construisons un unique
  !!{tableau} combiné ayant pour hypothèses
  $\sFmla{\True}{!B_1}$, \dots, $\sFmla{\True}{!B_n}$ ainsi que
  $\sFmla{\True}{!C_1}$, \dots, $\sFmla{\True}{!C_m}$, puis appliquons
  la règle \Cut{}. Nous obtenons deux branches : la première commence
  par $\sFmla{\True}{!A}$, la seconde par $\sFmla{\False}{!A}$.
  
  Sur la branche de gauche, greffons la partie du premier !!{tableau}
  située sous ses hypothèses. Chaque application de règle y demeure
  correcte, puisque toutes les hypothèses du premier !!{tableau}, y
  compris $\sFmla{\True}{!A}$, sont disponibles. Toutes les branches
  issues de $\sFmla{\True}{!A}$ se ferment donc.
  
  Sur la branche de droite, greffons la partie du second !!{tableau}
  située sous ses hypothèses, après avoir supprimé les résultats des
  applications de~$\TRule{\True}{\lnot}$ à
  $\sFmla{\True}{\lnot !A}$. La conclusion de
  $\TRule{\True}{\lnot}$ est
  $\sFmla{\False}{!A}$, qui reste disponible puisqu'elle est la
  conclusion droite de la règle \Cut{} dans le !!{tableau} combiné.

  Si une branche du second tableau se fermait parce qu'elle contenait
  l'hypothèse $\sFmla{\True}{\lnot !A}$, désormais absente du
  !!{tableau} combiné, ainsi que $\sFmla{\False}{\lnot !A}$, appliquons
  $\TRule{\False}{\lnot}$ à cette dernière formule. On obtient
  $\sFmla{\True}{!A}$; la branche correspondante du !!{tableau}
  combiné se ferme alors avec la conclusion droite de \Cut{},
  $\sFmla{\False}{!A}$. Si une branche du second !!{tableau} se fermait
  pour une autre raison, la branche correspondante du !!{tableau}
  combiné se ferme elle aussi, puisque toutes les !!{signed formula}s
  de l'ancienne branche autres que $\sFmla{\True}{\lnot !A}$ figurent
  encore sur la branche correspondante du !!{tableau} combiné.
  \end{proof}

\end{document}
